% Part: computability
% Chapter: computability-theory
% Section: universal-part-function

\documentclass[../../../include/open-logic-section]{subfiles}

\begin{document}

\olfileid{cmp}{thy}{uni}
\olsection{సార్వత్రిక పాక్షిక గణనీయ ప్రమేయం}

\begin{thm}
\ollabel{thm:univ-comp}
సార్వత్రిక పాక్షిక గణనీయ ప్రమేయం~$\fn{Un}(e,x)$ ఒకటి ఉంది. మరో
విధంగా చెప్పాలంటే, కింది ధర్మాలు గల ప్రమేయం $\fn{Un}(e,x)$ ఒకటి ఉంది:
\begin{enumerate}
\item $\fn{Un}(e,x)$ పాక్షిక గణనీయ ప్రమేయం.
\item $f(x)$ ఏదైనా పాక్షిక గణనీయ ప్రమేయమైతే, ప్రతి~$x$కూ
$f(x) \simeq \fn{Un}(e,x)$ అయ్యే సహజ సంఖ్య $e$ ఒకటి ఉంటుంది.
\end{enumerate}
\end{thm}

\begin{proof}
క్లీనీ నియత రూప సిద్ధాంతం (\olref[nfm]{thm:normal-form})లోని $U$,
$T$లను తీసుకుని, $\fn{Un}(e,x) \simeq
U(\umin{s}{T(e,x,s)})$గా ఉంచండి.
\end{proof}

\begin{explain}
పాక్షిక గణనీయ ప్రమేయాలకు మనకు ఒక కార్యసాధ్య లెక్కింపు ఉందని చెప్పే
కచ్చితమైన విధానమే ఇది. $f_e(x)=\fn{Un}(e,x)$గా నిర్వచించిన ప్రమేయాన్ని
$f_e$గా రాస్తే, $f_0$, $f_1$, $f_2$, \dots{} అనుక్రమంలో పాక్షిక
గణనీయ ప్రమేయాలన్నీ ఉంటాయి; పైగా $f_e(x)$ను $e$,~$x$లలో
``ఏకరీతిగా'' గణించవచ్చు. సరళత కోసం ఏకస్థాన ప్రమేయాలకు సార్వత్రికమైన
ద్విస్థాన ప్రమేయాన్ని వాడుతున్నాం; కానీ సంఖ్యల అనుక్రమాలను సంకేతీకరించి,
దీన్ని మరిన్ని ఆర్గుమెంట్లకు సులభంగా సాధారణీకరించవచ్చు. ఉదాహరణకు,
$f(x,y,z)$ ఒక $3$-స్థాన పాక్షిక పునరావృత్త ప్రమేయమైతే,
$g(x) \simeq f((x)_0,(x)_1,(x)_2)$ ఒక ఏకస్థాన పాక్షిక పునరావృత్త
ప్రమేయమని గమనించండి.
\end{explain}

\end{document}
